\section{Meta 评价与长期泛化}
\subsection{Meta-Judge 训练的细节}
Meta judge 的出现是因为 self-rewarding 模型在 few iterations 后会 plateau：Weston 形象地说 ``self-rewarding models that I showed earlier only had the top part of the pipeline; introducing meta judge builds the bottom part, lifting the plateau''（SRT 2756-2789）。Meta judge 的输入是 draft 与 verified response 对，输出是哪个更值得保留，训练目标不再依赖人工偏好，而是根据 Alpaca Eval 2 等 benchmark 的 win rate 进行强化。

\begin{figure}[H]
\centering
\includegraphics[width=0.77\textwidth]{frame-04-meta.jpg}
\caption{Meta Judge 章展示 draft、judge、verified 三步的分工。\protect\footnotemark}
\end{figure}
\footnotetext{视频画面时间区间：01:01:30--01:01:32。}

\subsection{Meta-Judge 数据与 SelfT Evaluators}
Weston 提到 ``SelfT evaluators 是我们去年 SelfT paper 做出来的''，该数据集把 verified response 的判断过程结构化：它包含 draft、verified、judge 三个阶段，而且 judge 的 label 来自桌面推理任务而非纯人工偏好（SRT 2902-2904）。利用 SelfT evaluators 训练出来的 judge 能更准确地区分出 ``真正提升 reasoning 质量'' 的回答，从而提升 win rate。

\begin{knowledgebox}{SelfT Evaluators 的设计逻辑}
SelfT Evaluators 训练用的是具有 chain-of-thought 的草稿和经过验证的答案，judge 只从两者之间挑出更可靠的那一个。这样训练出来的 evaluator 能以少量人监督或自监督数据，发挥比传统 reward 模型更强的分辨力。
\end{knowledgebox}

\subsection{推理可解释性与计算挑战}
Weston 认为 ``transformer itself does reasoning by all these steps''，但同一时间又承认网络本身并不透明；因此必须通过 meta judge 或可验证 reward 来追溯 reasoning trace（SRT 3152-3204）。他还强调 ``inference time compute for evaluation'' 是 self-improvement 的瓶颈，这也是为什么 verified response 数量不能随意增加，否则推理成本会翻倍（SRT 3236-3252）。

\subsection{本章小结}
Meta judge 与 SelfT evaluators 让 self-rewarding 模型具备了真正的判别能力：它们通过比较 draft 与 verified response，确保 reward signal 不被 hallucination 混淆，同时提醒我们推理计算必须可控。

